% 143-14(143쪽) — 곡선 $y=3x^2-6x$ 와 $x$축, 직선 $x=k$ 로 둘러싸인 두 도형(0~2 의 축 아래 부분 · 2~k 의 축 위 부분)에 색칠. 스캔 화질이 낮아 TikZ 로 다시 그림.
% 원문에 있는 것만 옮긴다: 곡선 · 직선 $x=k$ · 색칠한 두 도형 · 라벨 O, 2, k, x, y, $y=3x^2-6x$, $x=k$. 눈금값은 원문에 없으므로 적지 않았고, $x=k$ 는 답이 그림에서 읽히지 않도록 원문처럼 2 보다 조금 오른쪽(2 와 3 사이)에 두었다.
\begin{tikzpicture}[x=0.42cm, y=0.26cm, line width=0.6pt]
  \fill[black!12] plot[domain=0:2, samples=61, smooth] (\x,{3*(\x)*(\x)-6*(\x)}) -- cycle;
  \fill[black!12] plot[domain=2:2.6, samples=31, smooth] (\x,{3*(\x)*(\x)-6*(\x)}) -- (2.6,0) -- cycle;
  \draw[->, >=stealth, line width=0.7pt] (-1.5,0) -- (3.8,0) node[below=1pt, font=\small] {$x$};
  \draw[->, >=stealth, line width=0.7pt] (0,-4.3) -- (0,11.6) node[left=1pt, font=\small] {$y$};
  \draw[line width=0.5pt] (2.6,-3.4) -- (2.6,9.6);
  \draw[domain=-1.15:2.95, samples=121, smooth] plot (\x,{3*(\x)*(\x)-6*(\x)});
  \node[below left=-1pt and -4pt, font=\small] at (0,0) {$\pt{O}$};
  \node[below=1pt, font=\small] at (2,0) {$2$};
  \node[below right=0pt and -2pt, font=\small] at (2.6,0) {$k$};
  \node[font=\small] at (2.6,-4.5) {$x=k$};
  \node[anchor=west, font=\small] at (0.3,11.0) {$y=3x^2-6x$};
  \draw[->, >=stealth, line width=0.4pt] (4.0,10.2) to[out=-110, in=35] (3.05,8.2);
\end{tikzpicture}
